\chapter{ساختارها و شمارشی‌ها}%
\label{ch:structs}

تا اینجا با داده‌های ساده (عدد، رشته، منطقی) و مجموعه‌های هم‌نوع (آرایه، بردار)
کار کردیم. اما دنیای واقعی پیچیده‌تر است: یک «دانشجو» هم نام دارد، هم شماره، هم
معدل. چگونه این‌ها را با هم نگه داریم؟ پاسخ، \textbf{ساختار} است. در این فصل
ساختارها و خویشاوندشان، \textbf{شمارشی‌ها}، را می‌شناسیم.

\section{ساختار: ساختنِ نوعِ دلخواهِ خود}

ساختار به شما اجازه می‌دهد چند داده را زیرِ یک نام گرد آورید و یک \emph{نوعِ
تازه} بسازید. به این مثالِ واقعی نگاه کنید که یک «مستطیل» می‌سازد:

\begin{salam}
ساختار مستطیل:
    طول : اعشار۶۴ = 0.0
    عرض : اعشار۶۴ = 0.0
    تابع مساحت() : اعشار۶۴:
        بازگشت این.طول * این.عرض
    تمام
    تابع محیط() : اعشار۶۴:
        بازگشت 2.0 * (این.طول + این.عرض)
    تمام
تمام

تابع اصلی:
    متغیر r : مستطیل = مستطیل { طول = 4.0, عرض = 3.0 }
    چاپ("مساحت:", r.مساحت())
    چاپ("محیط:", r.محیط())
تمام
\end{salam}

این مثال چند مفهومِ تازه دارد که یکی‌یکی می‌شکافیم.

\subsection{میدان‌ها (fields)}

\begin{salam}
ساختار مستطیل:
    طول : اعشار۶۴ = 0.0
    عرض : اعشار۶۴ = 0.0
\end{salam}

\code{طول} و \code{عرض} \textbf{میدان}های ساختار هستند یعنی داده‌هایی که هر
مستطیل نگه می‌دارد. هر میدان نامی، نوعی و یک مقدارِ پیش‌فرضِ \emph{اختیاری} دارد.
مقدارِ پیش‌فرض (\code{= 0.0}) وقتی به کار می‌آید که هنگامِ ساختنِ مستطیل، مقداری
برای آن میدان ندهید.

مقدارِ پیش‌فرض اجباری نیست. میدانی که بدونِ \code{= مقدار} نوشته شود، یک میدانِ
\textbf{الزامی} است: هر نمونه باید آن را تعیین کند وگرنه کامپایلر خطا می‌دهد. به
این ترتیب داده‌ای را که ساختار بدونِ آن بی‌معناست الزامی می‌کنید و بقیه را به
پیش‌فرضِ معقول می‌سپارید:

\begin{salam}
ساختار کاربر:
    نام : رشته                 // الزامی بدونِ پیش‌فرض
    نقش : رشته = "عضو"         // اختیاری پیش‌فرض «عضو»
    فعال : منطقی = true
تمام

تابع اصلی:
    متغیر کا : کاربر = کاربر { نام = "سارا" }
    متغیر مدیر : کاربر = کاربر { نام = "مکس", نقش = "مدیر" }
    چاپ(کا.نام, کا.نقش, کا.فعال)
    چاپ(مدیر.نام, مدیر.نقش, مدیر.فعال)
تمام
\end{salam}

\begin{progout}
سارا عضو true
مکس مدیر true
\end{progout}

نوشتنِ \code{کاربر \{\}} اینجا خطای زمانِ کامپایل است، چون میدانِ الزامیِ
\code{نام} داده نشده. میدان‌های دارای پیش‌فرض را همیشه می‌توان ننوشت.

\subsection{ساختنِ یک نمونه}

\begin{salam}
متغیر r : مستطیل = مستطیل { طول = 4.0, عرض = 3.0 }
\end{salam}

این خط یک \textbf{نمونه} (instance) از مستطیل می‌سازد. نحوِ \code{مستطیل \{ ... \}}
یک مستطیلِ تازه با میدان‌های مشخص‌شده می‌سازد. حالا \code{r} یک مستطیلِ واقعی
با طولِ ۴ و عرضِ ۳ است.

\subsection{متدها: توابعِ درونِ ساختار}

\begin{salam}
تابع مساحت() : اعشار۶۴:
    بازگشت این.طول * این.عرض
تمام
\end{salam}

تابعی که \emph{درونِ} یک ساختار تعریف شود، \textbf{متد} نام دارد. متدها روی
داده‌های همان ساختار کار می‌کنند. واژهٔ کلیدیِ \code{این} (معادلِ \code{this})
به \emph{همان نمونه‌ای} اشاره می‌کند که متد رویش صدا زده شده. پس
\code{این.طول} یعنی «طولِ همین مستطیل».

متدها را با نقطه پس از نامِ نمونه صدا می‌زنیم: \code{r.مساحت()}. این دقیقاً
همان شیوه‌ای است که در فصلِ بردارها برای \code{.بیفزا()} و \code{.طول()}
دیدیم و حالا می‌دانید که آن‌ها هم متد بودند.

\section{تغییرِ داده‌ها از درونِ متد}

متدها فقط نمی‌خوانند؛ می‌توانند داده‌های ساختار را \emph{تغییر} هم بدهند. این
مثالِ واقعی یک «شمارنده» می‌سازد:

\begin{salam}
ساختار شمارنده:
    مقدار : صحیح = 0
    تابع افزایش():
        این.مقدار = این.مقدار + 1
    تمام
    تابع بگیر() : صحیح:
        بازگشت این.مقدار
    تمام
تمام

تابع اصلی:
    متغیر ش : شمارنده = شمارنده { مقدار = 5 }
    ش.افزایش()
    ش.افزایش()
    چاپ ش.بگیر()
تمام
\end{salam}

متدِ \code{افزایش} مقدارِ درونیِ شمارنده را یکی زیاد می‌کند. پس از دو بار صدا
زدن، مقدار از ۵ به ۷ می‌رسد و خروجی \code{7} می‌شود. این الگو نگه‌داشتنِ
داده در ساختار و تغییرِ آن از راهِ متدها بنیادِ برنامه‌نویسیِ شیءگراست.

\section{مثالِ کامل: حسابِ بانکی}

این مثالِ واقعی همهٔ آموخته‌ها را گرد می‌آورد و یک حسابِ بانکیِ ساده می‌سازد که
واریز و برداشت دارد:

\begin{salam}
ساختار حساب:
    موجودی : صحیح = 0
    تابع واریز(مبلغ : صحیح): این.موجودی = این.موجودی + مبلغ تمام
    تابع برداشت(مبلغ : صحیح):
        اگر مبلغ <= این.موجودی: این.موجودی = این.موجودی - مبلغ
        وگرنه:                   چاپ "موجودی کافی نیست"
        تمام
    تمام
تمام

تابع اصلی:
    متغیر ح : حساب = حساب { موجودی = 100 }
    ح.واریز(50)
    ح.برداشت(30)
    ح.برداشت(200)
    چاپ "موجودی =", ح.موجودی
تمام
\end{salam}

دقت کنید که متدِ \code{برداشت} پیش از کم‌کردن، بررسی می‌کند که موجودی کافی
باشد. این یک قاعدهٔ مهم در طراحیِ ساختارهاست: \emph{متدها از درستیِ داده‌ها
پاسداری می‌کنند.} برداشتِ ۲۰۰ از حسابی که فقط ۱۲۰ دارد، رد می‌شود و خروجیِ
نهایی \code{موجودی = 120} است.

\begin{tip}
به دو شیوهٔ دسترسی به میدان دقت کنید: از \emph{بیرونِ} ساختار با نامِ نمونه
(\code{ح.موجودی})، و از \emph{درونِ} متد با \code{این} (\code{این.موجودی}).
هر دو به همان داده اشاره می‌کنند.
\end{tip}

\section{عملگرها برای ساختارها}

در فصلِ ۳ به‌اختصار دیدیم که می‌توان عملگرها را برای ساختارها تعریف کرد. اکنون
که ساختارها را می‌شناسید، این مثالِ کامل‌تر روشن‌تر است:

\begin{salam}
ساختار بردار:
    ایکس : صحیح = 0
    وای : صحیح = 0
    تابع عملگر+(دیگری : بردار) : بردار:
        بازگشت بردار { ایکس = این.ایکس + دیگری.ایکس, وای = این.وای + دیگری.وای }
    تمام
    تابع عملگر==(دیگری : بردار) : منطقی:
        بازگشت این.ایکس == دیگری.ایکس && این.وای == دیگری.وای
    تمام
تمام

تابع اصلی:
    متغیر الف : بردار = بردار { ایکس = 1, وای = 2 }
    متغیر ب : بردار = بردار { ایکس = 3, وای = 4 }
    متغیر ج : بردار = الف + ب
    چاپ(ج.ایکس, ج.وای)
    چاپ(الف == ب)
    چاپ(ج == بردار { ایکس = 4, وای = 6 })
تمام
\end{salam}

با تعریفِ \code{عملگر+} و \code{عملگر==}، حالا \code{الف + ب} و
\code{الف == ب} روی بردارها معنا پیدا می‌کنند. خروجی به ترتیب: \code{4 6}،
\code{false} و \code{true}.

\begin{note}
سلام بسیاری عملگر را برای ساختارها می‌پذیرد: حسابی (\code{+ - * / \% **})،
مقایسه‌ای (\code{== != < > <= >=})، یکانی (\code{-} و \code{!}) و حتی اندیس‌گذاری
(\code{[]}). اگر \code{==} را تعریف کنید اما \code{!=} را نه، سلام به‌طورِ
خودکار \code{!=} را از روی \code{==} می‌سازد.
\end{note}

\section{شمارشی: مجموعه‌ای از نام‌های ثابت}

گاهی متغیری داریم که فقط چند مقدارِ مشخص می‌تواند بگیرد: روزهای هفته، جهت‌های
اصلی، رنگ‌های چراغِ راهنما. برای این کار، \textbf{شمارشی} (enum) را داریم. به
این مثالِ واقعی نگاه کنید:

\begin{salam}
شمارش رنگ:
    قرمز
    سبز = 5
    آبی
تمام

تابع اصلی:
    چاپ(رنگ.قرمز بعنوان صحیح)
    چاپ(رنگ.سبز بعنوان صحیح)
    چاپ(رنگ.آبی بعنوان صحیح)
تمام
\end{salam}

\code{شمارش رنگ} یک نوعِ تازه می‌سازد که سه مقدارِ ممکن دارد: \code{قرمز}،
\code{سبز} و \code{آبی}. هر عضوِ شمارشی در پسِ پرده یک عددِ صحیح است که از صفر
آغاز می‌شود. اما می‌توانید مقدارِ دلخواه هم بدهید: اینجا \code{سبز = 5} گذاشته‌ایم،
پس \code{قرمز} برابرِ ۰، \code{سبز} برابرِ ۵ و \code{آبی} برابرِ ۶ (یکی پس از
سبز) می‌شود. با \code{بعنوان صحیح} می‌توان عددِ پشتِ هر عضو را دید.

\subsection{شکلِ یک‌خطی و کاربردِ واقعی}

شمارشی‌ها را می‌توان فشرده‌تر هم نوشت. این مثالِ واقعیِ «روزهای هفته» شکلِ
یک‌خطی را نشان می‌دهد و شمارشی را در یک تابع به کار می‌برد:

\begin{salam}
شمارش روز: دوشنبه, سه‌شنبه, چهارشنبه, پنجشنبه, جمعه, شنبه, یکشنبه تمام

تابع آخرهفته است(ر : روز) : منطقی: بازگشت ر == روز.شنبه || ر == روز.یکشنبه تمام

تابع اصلی:
    ر : روز = روز.شنبه
    چاپ "شماره‌ی روز =", ر بعنوان صحیح
    اگر آخرهفته است(ر): چاپ "آخر هفته است"
    وگرنه:               چاپ "روز کاری است"
    تمام
تمام
\end{salam}

اینجا تابعِ \code{آخرهفته است} یک \code{روز} می‌گیرد و با مقایسهٔ آن با
\code{روز.شنبه} و \code{روز.یکشنبه}، تصمیم می‌گیرد. شمارشی‌ها کد را
\emph{خواناتر} و \emph{ایمن‌تر} می‌کنند: به‌جای عددهای جادوییِ بی‌معنا (۵، ۶)،
نام‌های گویا (\code{شنبه}، \code{یکشنبه}) به کار می‌بریم.

\begin{tip}
هر جا متغیری دارید که فقط چند حالتِ مشخص می‌گیرد وضعیتِ سفارش (در انتظار،
ارسال‌شده، تحویل‌شده)، نوعِ کاربر (مهمان، عضو، مدیر) یک شمارشی بسازید. این
کار از خطاهای فراوان جلوگیری می‌کند، چون کامپایلر نمی‌گذارد مقداری بیرون از
فهرست به آن متغیر بدهید.
\end{tip}

\section{خلاصهٔ فصل}

\begin{itemize}
  \item ساختار چند داده را زیرِ یک نام گرد می‌آورد و نوعی تازه می‌سازد.
  \item میدان‌ها داده‌های ساختارند و می‌توانند مقدارِ پیش‌فرضِ اختیاری داشته
        باشند؛ میدانِ بدونِ پیش‌فرض \emph{الزامی} است و باید در هر نمونه داده شود.
  \item متدها توابعِ درونِ ساختارند و با \code{این} به نمونهٔ جاری دسترسی
        دارند.
  \item نمونه را با نحوِ \code{نام \{ میدان = مقدار \}} می‌سازیم.
  \item عملگرها را می‌توان برای ساختارها تعریف کرد.
  \item شمارشی مجموعه‌ای از نام‌های ثابت است که کد را خوانا و ایمن می‌کند.
\end{itemize}

\section{تمرین‌ها}

\exercise ساختاری به نامِ \code{دایره} با میدانِ شعاع بسازید و متدهای
\code{مساحت} و \code{محیط} را برایش بنویسید (از $\pi \approx 3.1416$ استفاده
کنید).

\exercise ساختارِ \code{دانشجو} با میدان‌های نام، شماره و معدل بسازید و متدی
بنویسید که بگوید دانشجو مشروط است یا نه (معدلِ زیرِ ۱۲).

\exercise شمارشیِ \code{جهت} با اعضای شمال، جنوب، شرق و غرب بسازید و تابعی
بنویسید که جهتِ مقابلِ هر جهت را برگرداند.

\exercise به ساختارِ \code{بردار} متدِ \code{عملگر-} (تفریق) را نیز بیفزایید.

\exercise ساختارِ \code{حساب} را گسترش دهید تا تاریخچهٔ تراکنش‌ها را در یک بردار
نگه دارد.
